\documentclass[book.tex]{subfiles}
%\usepackage{bm}% bold math
%\usepackage{amsmath}

\begin{document}
% Julia Ling Paper can be covered alongside
% Cover Varun's recent paper too (divergence-free constraint)

% This chapter primarily motivates the use of ML in the context of turbulent flow within vehicle aerodynamics
\chapter{Convolutional Neural Ordinary Differential Equations for Fluid Modeling}

\label{chap:ml_fluid}
\noindent\rule{11cm}{0.4pt} \\
\textbf{Prerequisites:} \autoref{chap:qft_basics}  \\
\textbf{Difficulty Level:} ** \\
\noindent\rule{11cm}{0.4pt}
\vspace{5mm}

%In this chapter, we will motivate the use of machine learning methods in the context of fluid dynamics. We will specifically discuss a machine learning architecture that have shown promise in terms of predicting fluid flow fields close to the solutions of the Navier-Stokes equations. 

Turbulence modeling has traditionally taken a data-driven approach – closure models are developed, refined, and fit with empirical experimental or direct numerical simulation data via a-posteriori or a-priori tests. As machine learning methods mature, a multitude of new avenues are open to incorporating data into fluid flow prediction. While traditional computational fluid dynamics (CFD) techniques rely on known physics, empirical fitting, and mathematical derivations, machine learning techniques offer an alternative approach with enhanced flexibility and richer model representations. Many data-driven models have shown significant promise with their ability to improve or efficiently replace CFD methods. The potential for ML methods to achieve accurate, spatio-temporal fluid field surrogate prediction is compelling. 

Convolutional Neural ODEs can be used to approximate predict fluid flow fields \cite{shankar2020learning}. This chapter outlines this data-driven model for forecasting the velocity field of a turbulent flow and aims to make connections between the model and the physical intuitions of the system.

\subsection{Governing Equations of the System}

As an example, consider an idealized case of turbulent flow – homogeneous isotropic turbulence. For an incompressible fluid, there are two governing equations, a momentum balance and a mass balance, shown here in their non-dimensional form:
\begin{align}
\frac{\partial u}{\partial t}+(u\cdot\nabla)u&=\frac{1}{Re}\nabla^2u-\nabla p+f \\
\nabla\cdot u &=0
\label{eq:ns}
\end{align}
where $u$ is the velocity vector, defined $\{ u,v,w\}$, $p$ is the pressure and $Re$ is a Reynolds number, a non-dimensional flow parameter which broadly describes the intensity of inertial forces over viscous ones.  We can see from the momentum balance that the dynamics of the velocity field are both autonomous and a function of local spatial gradient information of the field. Given a spatially discrete field, then, we can construct an ODE system to describe the dynamics using a neural network. The convolutional neural ODE design stems from the fact that local spatial dynamics are critical for modeling time evolution.

\subsection{Convolution Operations Can Encode Spatial Derivatives}
In order to incorporate the spatial information (derivatives) necessary to accurately approximate the true function, we use a convolutional neural network to produce the dynamics. The memory costs of a convolutional network are drastically reduced as the number of parameters is much smaller than a fully-connected network. This is due to the network using a shared local filter in each layer that limits the message passing to a constrained receptive field around each output point. 

When we examine the true dynamics, we again see that they are only a function of local information and derivatives. Here, we provide the intuition behind this model structure by showing how fixed filter convolutions are functionally equivalent to derivative computation using finite differnce (FD) or finite volume (FV) techniques. Consider the 1D second order finite difference approximation for the second derivative of field variable $\phi$ at discrete location $i$:

\begin{align}
\frac{\partial^2 \phi_i}{\partial x^2} = \frac{\phi_{i+1} -2 \phi_{i} + \phi_{i-1}}{\Delta x^2} + O(\Delta x^2)
\end{align}

This operation is analogous to applying the convolutional kernel

\begin{align}
\frac{\partial^2 \phi}{\partial x^2} \approx \frac{1}{\Delta x^2} [1,-2,1] \quad ,
\end{align}

and this equivalency extends to 2D:

\begin{align*}
\frac{\partial \phi}{\partial x} &\approx \frac{1}{\Delta x} \begin{bmatrix}
0 & 0 & 0\\
-1 & 0 & 1\\
0 & 0 & 0
\end{bmatrix} \\
\frac{\partial \phi}{\partial y} &\approx \frac{1}{\Delta x} \begin{bmatrix}
0 & 1 & 0\\
0 & 0 & 0\\
0 & -1 & 0
\end{bmatrix}
\quad ,
\end{align*}


\noindent and 3D as well. Higher order difference schemes will naturally result in larger and larger kernel sizes. Although the convolution weights are not fixed, as they would be with a defined numerical discretization, we can treat the convolutional layers as a learned differencing scheme to account for higher-order behavior.

\section{Machine Learning Approach}

The model is provided with an initial snapshot of the flow. Continuous dynamics are forcasted with the model, and a series of temporal snapshots within a specified time window are saved and analyzed with regards to ground truth data.  We show here results using a series of four turbulent statistics – energy spectra, velocity probability density functions (PDFs), velocity autocorrelation, and turbulent kinetic energy (TKE) over time. The turbulent kinetic energy is defined as one half of the sum of the mean-square fluctuations of the velocity components:
\begin{align}
e &= \frac{1}{2}(\langle u^2\rangle+\langle v^2\rangle+\langle w^2\rangle) \quad 
\label{eq:tke}
\end{align}

\subsection{Convolutional AutoEncoder with NeuralODEs}

%A challenging aspect of this formulation is that the resultant system of ODEs is of the order $N^3$, corresponding to the cubic grid with side length $N$. 
Due to the combined multi-scale and non-linear behaviour of the Navier-Stokes equations, discretized numerical solutions to (\ref{eq:ns}) are subject to unfavorable scaling of computational complexity with respect to the Reynolds number.
To improve computation tractability, reduced order modeling lowers the dimensionality of dynamical systems. %Proper orthogonal decomposition, often used with fluid flows, takes a statistical approach, projecting the data onto a linear subspace comprised of dominant modes. The idea of a latent space representation is common in many fields, including computer vision.

Convolutional autoencoders have been used for tasks like image compression or feature extraction, both of which closely parallel the goal of a rich reduced representation of the turbulent flow field. We take advantage of this technique to reduce the snapshots into a much smaller latent space. The dynamics of the system can then be performed on the latent representation, easing computation. %Using a latent space in conjunction with Neural ODEs is not new and has been used successfully for modeling many other dynamical systems.

The encoder, denoted by $\mathcal{E}$, is a trainable network that transforms the input velocity field $u_0$.
\begin{align}
\mathcal{E}(u_0) = z_0
\end{align}
Compression is achieved with the use of strided convolutional layers in the encoder to reduce the spatial resolution in the latent space. This is balanced by an increase in the number of latent channels, up from the original 3 velocity channels. The compression ratio is defined as the size of the original data relative to the latent space. The initial condition in the latent space $z_0$ is used to solve the ODE parameterized by the network weights.
\begin{align}
\frac{d z}{d t}=g_\theta(z), \quad z(0) = z_0
\end{align}
Although this is a continuous function, and evaluated as such, the solution is saved at discrete time points corresponding to the training and test set.

The decoder, $\mathcal{D}$, is another trainable network with the same architecture as the encoder, and decodes the latent sequence back to the original velocity space.
\begin{align}
\mathcal{D}(\{z_0,z_1,\dotsc,z_n\}) = \{\hat{u}_0,\hat{u}_1,\dotsc,\hat{u}_n\}
\end{align}
This output sequence is used to train the networks using a normalized mean squared error (MSE) loss function. For this problem, we normalize the loss with the average energy in the flow.

The design of the Convolutional autoencoder + Neural ODE system is sketched in fig. \ref{fig:arch}.
%Here, the latent space dynamics $\frac{d z}{d t}$ are approximated using a 3 layer convolutional network with kernel sizes of 7. The encoder is a 2 layer convolutional network, with 1 layer of stride 2 to downsample the output. The decoder shares this architecture with convolutional transpose layers. We employ a compression ratio of 6, from halving the spatial resolution in each dimension and increasing the latent features to 4 ($2^3*3/4=6$). %Hyperparameters such as number of hidden layers and channels were selected from parameter sweeps and Bayesian optimization, although a consistent correlation with these features was not observed. All models were trained with the ADAM optimizer using a decaying learning rate from $10^{-3}$ - $10^{-5}$. Additionally, models were trained by progressively increasing the prediction horizon from $\tau=0.25$ to $\tau=1.0$.

\begin{figure}
  \centering
  % \includegraphics{figures/Differentiable Physics/fluid_ml/conv_node/arch_s_NODES_fluids.png}
    \includegraphics{figures/Differentiable Physics/fluid_ml/conv_node/neuralODE approximator.png}

  \caption{\small{Schematic of the overall model architecture. Initial conditions are given and encoded into a latent space. Augmented channels are concatenated and the dynamics are forecasted through the neural ODE approximator. The sequence is decoded and the divergence-free condition is enforced through the last layer.}
  }
  \label{fig:arch}

\end{figure}


\subsection{Dataset and Improved Boundary Constraints}


The model can be trained with high-fidelity solutions to the three-dimensional Navier-Stokes equations, defined by \ref{eq:ns}.
Solutions for a triply-periodic domain can be obtained by a Fourier pseudo-spectral method. %\textcolor{red}{TODO: What is a Fourier pseudo-spectral method}
The system should be forced at low-wavenumbers to keep the total energy in the system constant. The resulting data will then statistically stationary with approximately stationary energy spectra. %The training, validation and test datasets span 100 integral time scales $\tau$, with a temporal sampling rate of approximately $\tau/100$. The turbulent Reynolds number, $Re_T \equiv e^2/(\nu \epsilon)$ with $\epsilon$ the turbulent kinetic energy dissipation rate, is approximately 380. The solutions are discretized with 64 collocation points in each direction.%large eddy turnover times, the amount of time it takes for energy to traverse the cascade, with 100 individual time steps per large eddy time. 


In this approach, we can incorporate a few additional architectural elements to improve prediction and enforce physical constraints. As the dataset includes periodic boundary conditions, we impose this constraint by padding inputs to all convolutions circularly, as shown in Fig. \ref{fig:BC}. 

\begin{figure}
  \centering
  % \includegraphics[width=0.5\linewidth]{figures/Differentiable Physics/fluid_ml/conv_node/PBCs_NODES_fluids.png}
  \includegraphics[width=0.6\linewidth]{figures/Differentiable Physics/fluid_ml/conv_node/PBC_Nodes.png}

  \caption{\small{Inputs to convolutional operations are padded in a periodic fashion, preserving the boundary conditions specified in the problem dataset. Shown here is an example of padding a small 2D tensor. Analogously, this can be applied in 3D and with the appropriate amount of padding for the kernel size and valid outputs.}
  }
  \label{fig:BC}

\end{figure}

Within this approach, we also employ the augmented neural ODE approach to the latent dynamics to increase flexibility and ease computation of ODE. This is accomplished by concatenating channels of zeros to the output of the encoder, with the predicted results projected back into the original latent space. Lastly, we apply a spectral projection layer to the final model output to enforce the divergence-free field condition required for a constant density fluid. 

\subsection{Spectral Projection for Divergence-Free Constraint}

To add a divergence-free constraint, the input $u'$ is first transformed into Fourier space:

\begin{align}
    U'=\mathcal{F}(u') \quad ,
\end{align}

where the divergence operator can be applied linearly. 

With this linear constraint, we can construct a quadratic optimization problem to minimize the squared difference between the input and output in Fourier space: 

\begin{align}
    \min_{\hat{U}} \quad & \frac{1}{2}(U'-\hat{U})^T(U'-\hat{U}) = \\
    \min_{\hat{U}} \quad & \frac{1}{2}\hat{U}^T I\hat{U} -
    (U')^T\hat{U}\\
    \textrm{s.t.} \quad & A\hat{U} = 0 \quad ,
\end{align}

where $A$ is the linear divergence operator. The solution can be obtained directly:

\begin{align}
    \hat{U} = U'-\frac{k\cdot U'}{k\cdot k}k \quad ,
\end{align}

where $k$ are the wavenumbers from the Fourier analysis. The final output comes from taking the inverse Fourier transform of the optimum.

\begin{align}
    \hat{u}=\mathcal{F}^{-1}(\hat{U})
\end{align}

Each operation in the layer is fully differentiable and thus backpropagation through the layer with automatic differentiation software is straightforward. 


\subsection{Evaluation Criteria}

We can analyze the predictive capabilities of the model by examining a series of turbulent statistical metrics and comparing them to those of the DNS dataset. While standard statistical measures like RMSE are often employed in machine learning applications to judge performance of models, these values are limited in their physical insight and ultimately only offer a holistic qualitative measure of accuracy compared to the true data. Rigorous scientific analysis of machine learning models in this space requires a physical understanding of the results. Hence, we look at five flow statistics: %that have been well documented analytically and quantitatively which provide different perspectives on the flow and can afford insight into how the model behaves in accordance with the physical knowledge of the system.

%We look at 

\begin{enumerate}
\item Energy spectra: the energy spectrum decomposes the turbulent kinetic energy in the flow as a function of wavenumber, where large and small eddies correspond to low and high wavenumbers respectively. It is well-known that 3D turbulent flow exhibits a power law behavior in the inertial range of the energy spectrum in accordance with Kolmogorov’s 4/5-law. The spectra can tell us information about how the model performs at different scales in the flow. 
%\item Intermittency: \textcolor{red}{TODO}.
\item Q-R plots: real turbulent flows exhibit certain universal small-scale structures, such as a preference for vorticity alignment with the intermediate strain-rate eigenvalue. The flow topology can be described using the Q-R plane, which represent the second and third velocity gradient tensor invariants respectively. It has been found that turbulent flows show a characteristic asymmetrical teardrop shape in the joint Q-R pdf, which is not apparent in a random field. %We additionally perform coarse Q-R tests, which provide a description of this flow topology at varying length scales.
\end{enumerate} 

%We also look at two temporally varying statistics. 
\begin{enumerate}
  \setcounter{enumi}{3}
    \item Velocity autocorrelation: the autocorrelation function describes the similarity of a dynamical system’s state to itself over time. In turbulent flows, the function decays rapidly over time and the integral converges to a value defined as the integral timescale of the process. Given the stationarity imposed by the dataset, we expect to see additional power law behavior in the correlation function within this integral timescale window.
    \item The turbulent kinetic energy in the flow over time, is defined in eq. \ref{eq:tke} and also given by the integral of the energy spectra. For this stationary process, we expect to see constant energy in the system.
\end{enumerate}

Fig. \ref{fig:qual} provides a qualitative visual comparison of the output of the convolutioonal neural ODE with the DNS ground truth midway through the prediction window at $\tau=0.5$. %Generally, large-scale structures appear to be preserved, while smaller scale fluctuations are filtered out by the model. This behavior is justified in looking at the energy spectra in fig. \ref{fig:spec}. In the low-wavenumber region, the energy in the flow is well-preserved, even over the integral timescale prediction horizon. At around $k=10$, we begin to see some deviation from the DNS data, where the model underpredicts the energy at those scales. This indicates that the encoder-neural-ODE architecture is able to capture the large-scale behavior but loses fidelity in approximating the small-scale energy. The intermittency plots in fig. \ref{fig:pdf} support this as well, with good agreement in the center and discrepancies at the tails of the pdf, which correspond to the small scales. We note that for many engineering applications, resolving these largest scales is often sufficient to realize dynamics of interest for engineering or geophysical problems.

\begin{figure}
  \centering
  \includegraphics[width=\linewidth]{figures/Differentiable Physics/fluid_ml/conv_node/qual__fluids_NODEs.png}

  \caption{Qualitative comparison between the stochastic fluid flow fields generated from DNS and the neural network approach presented in this chapter.}
  \label{fig:qual}
  
\end{figure}

\begin{marginfigure}
  \centering
  \includegraphics[width=\linewidth]{figures/Differentiable Physics/fluid_ml/conv_node/energy_spectrum_fluids_NODEs.pdf}

  \caption{\small{Both plots show the energy of the flow as a function of wavenumber. Left: Two snapshots corresponding to the initial and ending time steps are shown. The vertical line represents the compression ratio of 6 defined by the latent space size. Right: The energy spectra averaged over the entire time window is shown. Low wavenumbers show better agreement with DNS, which are sufficient for many practical flows.}}
  \label{fig:spec}

\end{marginfigure}

%The statistical tests we have demonstrated so far have examined the flow either at singular temporal snapshots or in a time-averaged fashion. It is also important to consider how the predictions vary in time, as part of the value proposition of the neural ODE framework is its ability to model the continuous field dynamics. To this end, we present 
The velocity autocorrelation, the turbulent kinetic energy, and the normalized MSE over the prediction horizon are presented in fig. \ref{fig:temporal}. We again note the good agreement and power-law behavior in the autocorrelation that matches the DNS dynamics. Importantly, we find that the network is adept at maintaining constant energy in the system over time. %Though the model consistently underpredicts the true energy, this is in line with findings from the energy spectra, and we point out that the small-scale energy losses only amount to less than 7\% of the total TKE. 

%Additionally, fig. \ref{fig:temporal} shows the magnitude of divergence in the velocity field over one integral timescale for two models. The impact of the spectral projection step is significant in reducing the output divergence to magnitudes comparable to the true field. The formulation of the layer means inputs and outputs do not vary wildly, so training dynamics are not affected by using the layer, and MSE and other evaluation metrics are comparable between the models.



\begin{marginfigure}
  \centering
  \includegraphics[width=\linewidth]{figures/Differentiable Physics/fluid_ml/conv_node/intermittency_fluids_NODEs.pdf}

  \caption{\small{Distributions of the velocity gradients in the flow at initial and end temporal snapshots. Predictions exhibit tighter distributions, but they are consistent over time.}}
  \label{fig:pdf}

\end{marginfigure}

%\begin{figure}
%
%  \centering
%
%  \includegraphics[width=\linewidth]{figures/Differentiable Physics/fluid_ml/conv_node/temporal_fluids_NODEs.pdf}
%
%  \caption{\small{The model uses a divergence-constraining layer to induce physical solution fields. The magnitude of the divergence is plotted as a function of time for three cases: the true divergence, the prediction from the model without any constraint layer, and the prediction from the same model with the divergence constraint. The difference between the predictions is significant, and it is clear the constrained model performs close to the expected output.}} 
%  \label{fig:temporal}
%
%\end{figure}

\section{Summary}

In this chapter, we have looked at the use of convolutional neural ODEs for forecasting complex nonlinear spatiotemporal dynamics through purely data-driven methods. This methodology uses a trainable encoder-decoder pair to transform the field data into a latent space that retains the spatial nature of the data and solve the transient problem on this space using standard ODE methods. We find that this approach is able to capture important details of the flow, namely the large-scale dynamics, and produces stationary sequences in line with the physics of the problem (Navier-Stokes Equations).

%\textcolor{red}{TODO: Include a physika code from the core implementation.}


\end{document}